\documentclass[11pt,a4paper]{article}

\usepackage{amsmath,amssymb}
\usepackage[margin=2cm]{geometry}
\usepackage{enumitem}

\setlength{\parindent}{0pt}
\setlength{\parskip}{4pt}
\setlength{\emergencystretch}{1.5em}
\setlist[enumerate,1]{label=\arabic*.,leftmargin=*,itemsep=0.7em}
\setlist[enumerate,2]{label=\alph*.,leftmargin=*,itemsep=0.3em,topsep=0.3em}
\newcommand{\vect}[1]{\overrightarrow{#1}}

\begin{document}

\begin{center}
    {\large\bfseries QUIZ FISIKA MATEMATIKA 1}\\
    {\bfseries SEMESTER GASAL 2026/2027}\\
    {\bfseries INSTITUT TEKNOLOGI SEPULUH NOPEMBER}
\end{center}
\begin{tabular}{@{}ll@{}}
    Mata Kuliah  & : Fisika Matematika 1 \\
    Hari/Tanggal & : Senin, 12 Oktober 2026 (12.30--14.10 WIB) \\
    Waktu        & : 100 menit (\textbf{Kejujuran di atas segalanya})
\end{tabular}

\smallskip
\hrule
\medskip
\textbf{Selesaikan semua soal secara runut dan berurutan!}

\begin{enumerate}
    \item Diketahui $A=(1,-1,2)$, $B=(2,1,-1)$, dan $C=(1,0,-3)$
    merupakan titik-titik di $\mathbb{R}^3$. Tentukan:
    \begin{enumerate}
        \item Luas segitiga $ABC$.
        \item Proyeksi orthogonal ruas garis $AB$ terhadap ruas garis yang
        tegak lurus terhadap ruas garis $AC$ dan $BC$. Tentukan normal
        bidang yang dibentuk oleh vektor $v$ dan $w$.
    \end{enumerate}

    \item Tentukan persamaan bidang yang melalui titik $u=(3,2,1)$,
    $v=(2,1,-1)$, dan $w=(-1,3,2)$.

    \item Tentukan nilai-nilai yang sesuai untuk $x_1$, $x_2$, $x_3$, dan
    $x_4$ untuk persamaan matriks berikut:
    \begin{enumerate}
        \item
        \[
        \begin{pmatrix}
            x_1+x_2+x_3 & 0 & 1 \\
            2 & 3 & 2x_1+4x_2-3x_3 \\
            4 & 3x_1+6x_2-5x_3 & 5
        \end{pmatrix}
        =
        \begin{pmatrix}
            9 & 0 & 1 \\
            2 & 3 & 1 \\
            4 & 0 & 5
        \end{pmatrix}.
        \]
        \item
        \[
        \begin{pmatrix}
            x_1-x_2 & x_2+x_3 \\
            3x_4+x_3 & 2x_1-4x_4
        \end{pmatrix}
        =
        \begin{pmatrix}
            8 & 1 \\
            7 & 6
        \end{pmatrix}.
        \]
    \end{enumerate}

    \item Dengan menggunakan pendekatan reduksi baris (metode
    Gauss--Jordan), tentukan balikan (invers) dari matriks $A$ berikut:
    \[
    A=\begin{pmatrix}
        1 & 2 & 3 \\
        2 & 5 & 3 \\
        1 & 0 & 8
    \end{pmatrix}.
    \]

    \item Hitunglah jarak terpendek dari titik $P=(3,2,3)$ ke bidang yang
    diberikan oleh $2x+y+2z=2$, dan temukan titik $Q$ pada bidang tersebut
    yang paling dekat dengan $P$.
\end{enumerate}

\newpage

\begin{center}
    {\Large\bfseries Pembahasan}
\end{center}

\begin{enumerate}
\item 
Vektor-vektor yang menghubungkan ketiga titik adalah
\begin{align*}
    \vect{AB}&=B-A=(1,2,-3),\\
    \vect{AC}&=C-A=(0,1,-5),\\
    \vect{BC}&=C-B=(-1,-1,-2).
\end{align*}

\begin{enumerate}
\item Luas segitiga adalah setengah luas jajargenjang yang dibentuk oleh
$\vect{AB}$ dan $\vect{AC}$:
\[
    L=\frac12\left\|\vect{AB}\times\vect{AC}\right\|.
\]
Hasil perkalian silang kedua vektor tersebut ialah
\[
    \vect{AB}\times\vect{AC}
    =\begin{vmatrix}
        \mathbf{i} & \mathbf{j} & \mathbf{k} \\
        1 & 2 & -3 \\
        0 & 1 & -5
    \end{vmatrix}
    =(-7,5,1).
\]
Dengan demikian,
\[
    L=\frac12\sqrt{(-7)^2+5^2+1^2}
     =\frac{\sqrt{75}}2
     =\boxed{\frac{5\sqrt{3}}2\text{ satuan luas}}.
\]

\item \textit{Catatan:} Vektor $v$ dan $w$ pada kalimat terakhir soal 1(b)
belum didefinisikan. Sesuai konteks soal, pada pembahasan ini diambil
$v=\vect{AC}$ dan $w=\vect{BC}$.

Arah yang tegak lurus terhadap $\vect{AC}$ dan $\vect{BC}$ diperoleh dari
perkalian silang:
\[
    \mathbf{n}=v\times w
    =\begin{vmatrix}
        \mathbf{i} & \mathbf{j} & \mathbf{k} \\
        0 & 1 & -5 \\
        -1 & -1 & -2
    \end{vmatrix}
    =(-7,5,1).
\]
Proyeksi vektor $\vect{AB}$ pada arah $\mathbf{n}$ adalah
\begin{align*}
    \operatorname{proj}_{\mathbf{n}}\vect{AB}
    &=\frac{\vect{AB}\cdot\mathbf{n}}{\mathbf{n}\cdot\mathbf{n}}
      \,\mathbf{n}\\
    &=\frac{(1)(-7)+(2)(5)+(-3)(1)}{75}(-7,5,1)\\
    &=\boxed{(0,0,0)}.
\end{align*}
Jadi panjang proyeksi ruas garis $AB$ pada arah tersebut adalah
$\boxed{0}$. Hal ini juga mengikuti dari
$\vect{AB}=\vect{AC}-\vect{BC}$: vektor $\vect{AB}$ berada pada bidang
yang tegak lurus terhadap $\mathbf{n}$.

Vektor normal bidang yang dibentuk oleh $v$ dan $w$ dapat dipilih sebagai
\[
    \boxed{\mathbf{n}=(-7,5,1)},
    \qquad
    \boxed{\widehat{\mathbf{n}}=\frac{1}{5\sqrt{3}}(-7,5,1)}
\]
untuk normal satuannya. Kelipatan skalar tak nol dari $\mathbf{n}$ juga
merupakan vektor normal, termasuk arah yang berlawanan.

\end{enumerate}

\item 
Dua vektor arah pada bidang diperoleh dengan mengurangkan titik $u$:
\[
    \mathbf{a}=v-u=(-1,-1,-2),\qquad
    \mathbf{b}=w-u=(-4,1,1).
\]
Vektor normalnya adalah
\[
    \mathbf{n}=\mathbf{a}\times\mathbf{b}
    =\begin{vmatrix}
        \mathbf{i} & \mathbf{j} & \mathbf{k} \\
        -1 & -1 & -2 \\
        -4 & 1 & 1
    \end{vmatrix}
    =(1,9,-5).
\]
Karena bidang melalui $u=(3,2,1)$, persamaannya memenuhi
\begin{align*}
    \mathbf{n}\cdot\bigl((x,y,z)-u\bigr)&=0,\\
    (x-3)+9(y-2)-5(z-1)&=0.
\end{align*}
Jadi persamaan bidangnya adalah
\[
    \boxed{x+9y-5z=16}.
\]
Ketiga titik memenuhi persamaan ini:
$3+18-5=16$, $2+9+5=16$, dan $-1+27-10=16$.

\item 
Dua matriks sama apabila setiap elemen yang seletak bernilai sama.
Nilai variabel pada bagian (a) dan (b) dicari secara terpisah.

\begin{enumerate}
\item
Dari elemen yang memuat variabel, diperoleh sistem
\begin{align}
    x_1+x_2+x_3&=9,\label{eq:3a-1}\\
    2x_1+4x_2-3x_3&=1,\label{eq:3a-2}\\
    3x_1+6x_2-5x_3&=0.\label{eq:3a-3}
\end{align}
Ambil tiga kali persamaan \eqref{eq:3a-2} dikurangi dua kali persamaan
\eqref{eq:3a-3}:
\[
    (6x_1+12x_2-9x_3)-(6x_1+12x_2-10x_3)=3-0,
    \qquad x_3=3.
\]
Substitusi $x_3=3$ ke persamaan \eqref{eq:3a-1} dan \eqref{eq:3a-2}
memberikan
\[
    x_1+x_2=6,\qquad x_1+2x_2=5.
\]
Selisih kedua persamaan menghasilkan $x_2=-1$, sehingga $x_1=7$.
Maka
\[
    \boxed{x_1=7,\qquad x_2=-1,\qquad x_3=3}.
\]
Pemeriksaan: $7-1+3=9$, $14-4-9=1$, dan $21-6-15=0$.
Variabel $x_4$ tidak muncul pada bagian (a), sehingga tidak dibatasi
oleh persamaan matriks ini.

\item 
Kesamaan elemen memberikan
\[
    \begin{cases}
        x_1-x_2=8,\\
        x_2+x_3=1,\\
        3x_4+x_3=7,\\
        2x_1-4x_4=6.
    \end{cases}
\]
Dari dua persamaan pertama,
$x_1=x_2+8$ dan $x_3=1-x_2$.
Substitusi ke persamaan ketiga dan keempat menghasilkan
\begin{align*}
    3x_4+1-x_2=7&\quad\Longrightarrow\quad x_2=3x_4-6,\\
    2(x_2+8)-4x_4=6&\quad\Longrightarrow\quad x_2=2x_4-5.
\end{align*}
Samakan kedua bentuk $x_2$:
\[
    3x_4-6=2x_4-5\quad\Longrightarrow\quad x_4=1.
\]
Selanjutnya $x_2=-3$, $x_1=5$, dan $x_3=4$. Jadi
\[
    \boxed{x_1=5,\qquad x_2=-3,\qquad x_3=4,\qquad x_4=1}.
\]
Pemeriksaan dengan substitusi memberikan
\[
    \begin{pmatrix}
        5-(-3) & -3+4 \\
        3(1)+4 & 2(5)-4(1)
    \end{pmatrix}
    =\begin{pmatrix}8&1\\7&6\end{pmatrix}.
\]

\end{enumerate}

\item
\begin{align*}
    [A\mid I_3]&=
    \left(\begin{array}{ccc|ccc}
        1&2&3&1&0&0\\
        2&5&3&0&1&0\\
        1&0&8&0&0&1
    \end{array}\right)
    \xrightarrow{-2B_1+B_2}
    \left(\begin{array}{ccc|ccc}
        1&2&3&1&0&0\\
        0&1&-3&-2&1&0\\
        1&0&8&0&0&1
    \end{array}\right)
    \\[1em]
    &\xrightarrow{-B_1+B_3}
    \left(\begin{array}{ccc|ccc}
        1&2&3&1&0&0\\
        0&1&-3&-2&1&0\\
        0&-2&5&-1&0&1
    \end{array}\right)
    \xrightarrow{2B_2+B_3}
    \left(\begin{array}{ccc|ccc}
        1&2&3&1&0&0\\
        0&1&-3&-2&1&0\\
        0&0&-1&-5&2&1
    \end{array}\right)
    \\[1em]
    &\xrightarrow{-B_3}
    \left(\begin{array}{ccc|ccc}
        1&2&3&1&0&0\\
        0&1&-3&-2&1&0\\
        0&0&1&5&-2&-1
    \end{array}\right)
    \xrightarrow{-3B_3+B_1}
    \left(\begin{array}{ccc|ccc}
        1&2&0&-14&6&3\\
        0&1&-3&-2&1&0\\
        0&0&1&5&-2&-1
    \end{array}\right)
    \\[1em]
    &\xrightarrow{3B_3+B_2}
    \left(\begin{array}{ccc|ccc}
        1&2&0&-14&6&3\\
        0&1&0&13&-5&-3\\
        0&0&1&5&-2&-1
    \end{array}\right)
    \xrightarrow{-2B_2+B_1}
    \left(\begin{array}{ccc|ccc}
        1&0&0&-40&16&9\\
        0&1&0&13&-5&-3\\
        0&0&1&5&-2&-1
    \end{array}\right).
\end{align*}
Jadi
\[
    \boxed{A^{-1}=\begin{pmatrix}
        -40&16&9\\
        13&-5&-3\\
        5&-2&-1
    \end{pmatrix}}.
\]
Sebagai pemeriksaan,
\[
    AA^{-1}=
    \begin{pmatrix}1&2&3\\2&5&3\\1&0&8\end{pmatrix}
    \begin{pmatrix}-40&16&9\\13&-5&-3\\5&-2&-1\end{pmatrix}
    =\begin{pmatrix}1&0&0\\0&1&0\\0&0&1\end{pmatrix}.
\]

\item 
Bidang $2x+y+2z=2$ mempunyai vektor normal $\mathbf{n}=(2,1,2)$
dengan $\|\mathbf{n}\|=3$.

\begin{enumerate}
\item Jarak titik $(x_0,y_0,z_0)$ ke bidang $ax+by+cz=d_0$ adalah
\[
    D=\frac{|ax_0+by_0+cz_0-d_0|}{\sqrt{a^2+b^2+c^2}}.
\]
Untuk $P=(3,2,3)$, diperoleh
\[
    D=\frac{|2(3)+2+2(3)-2|}{\sqrt{2^2+1^2+2^2}}
     =\frac{12}{3}
     =\boxed{4\text{ satuan panjang}}.
\]

\item 
Titik terdekat adalah kaki garis tegak lurus dari $P$ ke bidang.
Karena garis ini searah dengan normal bidang, tuliskan
\[
    Q=P-t\mathbf{n}=(3-2t,\,2-t,\,3-2t).
\]
Syarat $Q$ berada pada bidang memberikan
\begin{align*}
    2(3-2t)+(2-t)+2(3-2t)&=2,\\
    14-9t&=2,\\
    t&=\frac43.
\end{align*}
Dengan demikian,
\[
    Q=\left(3-\frac83,\,2-\frac43,\,3-\frac83\right)
     =\boxed{\left(\frac13,\frac23,\frac13\right)}.
\]
Pemeriksaan:
\[
    2\left(\frac13\right)+\frac23+2\left(\frac13\right)=2,
    \qquad
    \|P-Q\|=\sqrt{\left(\frac83\right)^2+
                       \left(\frac43\right)^2+
                       \left(\frac83\right)^2}=4.
\]
Vektor $P-Q=\tfrac43(2,1,2)$ sejajar dengan normal bidang, sehingga
jarak tersebut merupakan jarak terpendek.

\end{enumerate}
\end{enumerate}

\end{document}
